\documentclass[10pt,a4paper]{article}
\usepackage[T1]{fontenc}
\usepackage[utf8]{inputenc}
\usepackage{lmodern}
\usepackage{amsmath,amssymb,amsthm,mathtools}
\usepackage[margin=0.88in,top=0.92in,bottom=0.88in]{geometry}
\usepackage{microtype}
\linespread{0.965}
\usepackage[dvipsnames]{xcolor}
\usepackage{titlesec}
\usepackage{enumitem}
\usepackage{fancyhdr}
\usepackage[hidelinks]{hyperref}

\definecolor{ink}{HTML}{17293E}
\definecolor{accent}{HTML}{275D8C}
\definecolor{muted}{HTML}{536477}
\color{ink}
\titleformat{\section}{\large\bfseries\color{accent}}{\thesection.}{0.55em}{}
\titleformat{\subsection}{\normalsize\bfseries\color{ink}}{\thesubsection}{0.55em}{}
\titlespacing*{\section}{0pt}{0.78em}{0.22em}
\titlespacing*{\subsection}{0pt}{0.50em}{0.12em}
\setlength{\parindent}{0pt}
\setlength{\parskip}{1.8pt}
\setlist[itemize]{leftmargin=1.35em,topsep=2pt,itemsep=1pt,parsep=0pt}
\setlist[enumerate]{leftmargin=1.5em,topsep=2pt,itemsep=1pt,parsep=0pt}
\setlength{\abovedisplayskip}{4.3pt}
\setlength{\belowdisplayskip}{4.3pt}
\setlength{\abovedisplayshortskip}{3pt}
\setlength{\belowdisplayshortskip}{3pt}
\setlength{\headheight}{20pt}
\pagestyle{fancy}
\fancyhf{}
\fancyhead[L]{\footnotesize\color{muted} STAGE 1 \quad / \quad MODULE 4}
\fancyhead[R]{\footnotesize\color{muted} MEASURE THEORY}
\fancyfoot[C]{\footnotesize\color{muted}\thepage}
\renewcommand{\headrulewidth}{0.3pt}
\renewcommand{\headrule}{\hbox to\headwidth{\color{muted!35}\leaders\hrule height \headrulewidth\hfill}}
\newcommand{\R}{\mathbb R}
\newcommand{\Q}{\mathbb Q}
\newcommand{\N}{\mathbb N}
\newcommand{\B}{\mathcal B}
\newcommand{\Lcal}{\mathcal L}
\newcommand{\key}[1]{\par\smallskip\textcolor{accent}{\textbf{Key point.}}\ #1\par}

\begin{document}
\begin{center}
  {\LARGE\bfseries Module 4: Lebesgue Measure}\par
  \vspace{5pt}
  {\small\color{muted}Measure axioms, continuity, outer measure, Carath\'eodory's criterion, and Lebesgue measurability}\par
  \vspace{7pt}
  {\color{accent}\rule{0.88\linewidth}{0.7pt}}
\end{center}
\vspace{-4pt}

\section{Measures}
\subsection{Definition and classical examples}
On a measurable space $(\Omega,\mathcal F)$, a \emph{measure} is a function
$\mu:\mathcal F\to[0,\infty]$ satisfying
\[
\boxed{\mu(\varnothing)=0,\qquad
\mu\!\left(\bigsqcup_{n=1}^{\infty}A_n\right)
=\sum_{n=1}^{\infty}\mu(A_n)\quad(A_n\in\mathcal F).}
\]
Here $\bigsqcup$ denotes disjoint union. The triple
$(\Omega,\mathcal F,\mu)$ is a \emph{measure space}; $+\infty$ is allowed.
\begin{itemize}
\item \textbf{Counting measure:} on $(\Omega,\mathcal P(\Omega))$,
$\mu(A)=|A|$ if $A$ is finite, and $+\infty$ otherwise.
\item \textbf{Dirac measure:} for $x_0\in\Omega$,
$\delta_{x_0}(A)=\mathbf1_A(x_0)$ on any $\mathcal F$.
\item \textbf{Probability measure:} a measure $\mathbb P$ with
$\mathbb P(\Omega)=1$.
\end{itemize}

\subsection{Monotonicity and countable subadditivity}
\textbf{Monotonicity.} If $A\subseteq B$ are measurable, then
\[
\mu(B)=\mu(A)+\mu(B\setminus A)\ge\mu(A).
\]
\textbf{Countable subadditivity.} For arbitrary measurable $A_n$, disjointness
is not needed for
\[
\boxed{\mu\!\left(\bigcup_{n=1}^{\infty}A_n\right)
\le\sum_{n=1}^{\infty}\mu(A_n).}
\]
\emph{Proof idea.} Disjointify via $B_1=A_1$ and
$B_n=A_n\setminus\bigcup_{k<n}A_k$: then
$\bigcup_nB_n=\bigcup_nA_n$ and $B_n\subseteq A_n$.

\subsection{Finite and sigma-finite measures}
A measure is \emph{finite} if $\mu(\Omega)<\infty$. It is
\emph{sigma-finite} if $\Omega=\bigcup_{n\ge1}E_n$ for measurable $E_n$
with $\mu(E_n)<\infty$. Every finite measure is sigma-finite, but the
converse fails. Lebesgue length on $\R$ is sigma-finite, since
\[
\R=\bigcup_{n\ge1}[-n,n],\qquad \lambda([-n,n])=2n<\infty,
\]
although $\lambda(\R)=\infty$. Counting measure on an uncountable set
is not sigma-finite.

\key{A measure is countably \emph{additive} on disjoint sets and countably
\emph{subadditive} on arbitrary measurable sets. These are distinct assertions.}

\section{Continuity of Measures}
\subsection{Continuity from below and from above}
For increasing measurable sets $A_n\uparrow A$, meaning
$A_n\subseteq A_{n+1}$ and $A=\bigcup_nA_n$, one has
\[
\boxed{\mu(A_n)\uparrow\mu(A),\qquad
\mu\!\left(\bigcup_nA_n\right)=\lim_{n\to\infty}\mu(A_n).}
\]
\emph{Proof.} Put $B_1=A_1$, $B_n=A_n\setminus A_{n-1}$ ($n\ge2$).
These are disjoint and $A_m=\bigsqcup_{n\le m}B_n$,
$A=\bigsqcup_{n\ge1}B_n$; apply countable additivity.

For decreasing sets $A_n\downarrow A$, meaning
$A_n\supseteq A_{n+1}$ and $A=\bigcap_nA_n$, we have
\[
\boxed{\mu(A_1)<\infty\quad\Longrightarrow\quad
\mu(A_n)\downarrow\mu(A).}
\]
\emph{Proof.} Set $B_n=A_1\setminus A_n\uparrow A_1\setminus A$.
Continuity from below applies to $B_n$. Since $\mu(A_1)<\infty$,
subtraction in $\mu(A_n)=\mu(A_1)-\mu(B_n)$ is legitimate, and taking
limits gives $\mu(A)$.

\emph{Why the finiteness condition matters.} For Lebesgue measure, let
$A_n=[n,\infty)$. Then $A_n\downarrow\varnothing$ but
\[
\boxed{\mu(A_n)=\infty\text{ for every }n,\qquad
\mu\!\left(\bigcap_nA_n\right)=0.}
\]
Every point eventually leaves these infinite-length tails; subtraction
would require the undefined $\infty-\infty$.

\subsection{Connection with set limsup and liminf}
For arbitrary sets, recall the Module 1 definitions
\[
\limsup_n A_n=\bigcap_N\bigcup_{n\ge N}A_n,
\qquad
\liminf_n A_n=\bigcup_N\bigcap_{n\ge N}A_n.
\]
For monotone sequences, $\liminf A_n=\limsup A_n$ equals
$\bigcup_nA_n$ (increasing) or $\bigcap_nA_n$ (decreasing).
Measure continuity preserves these limits, with finiteness from above;
it is an \emph{order-sequential}, not unrestricted, continuity.

\emph{Counterexample without monotonicity.} For $A_n=[n,n+1]$,
$\liminf A_n=\limsup A_n=\varnothing$, yet $\lambda(A_n)=1$ for all $n$.

\key{Continuity means compatibility with limits. Countable additivity
implies continuity from below; finite mass prevents a loss of mass
``at infinity'' when sets decrease.}

\section{Outer Measure}
\subsection{Constructing Lebesgue outer measure}
How can interval length be extended to general sets? Cover the target
set by countably many open intervals and minimize the total covering
length. For every $A\subseteq\R$, define
\[
\boxed{\lambda^*(A)=\inf\left\{\sum_{n=1}^{\infty}|I_n|:
A\subseteq\bigcup_{n=1}^{\infty}I_n,\ 
I_n\text{ open intervals}\right\}.}
\]
The infimum need not be attained; $\lambda^*$ is defined on
\emph{all} subsets of $\R$. An \emph{outer measure}
$\mu^*:\mathcal P(\Omega)\to[0,\infty]$ satisfies
\[
\boxed{\mu^*(\varnothing)=0,\quad
A\subseteq B\Rightarrow\mu^*(A)\le\mu^*(B),\quad
\mu^*\!\left(\bigcup_n A_n\right)\le\sum_n\mu^*(A_n).}
\]
For $\lambda^*$, monotonicity follows by reusing covers;
subadditivity follows by combining near-optimal covers.

\subsection{Basic examples}
\textbf{Intervals.} For $a<b$, $\lambda^*([a,b])=b-a$. An interval
$(a-\varepsilon,b+\varepsilon)$ gives the upper bound $b-a+2\varepsilon$.
For the lower bound, compactness gives a finite subcover; a chain
of overlapping intervals from $a$ to $b$ shows its lengths sum to at least
$b-a$. The original covering sum is no smaller.

\textbf{Countable sets.} If $A=\{x_1,x_2,\ldots\}$, cover $x_n$ by an
open interval of length $\varepsilon/2^n$. The total length is
$\varepsilon$, so $\lambda^*(A)=0$. In particular,
\[
\boxed{\lambda^*(\Q\cap[0,1])=0,\qquad
\overline{\Q\cap[0,1]}=[0,1].}
\]
Thus topological density does not imply positive measure.

\subsection{Why outer measure is not yet a measure}
Outer measure guarantees only \emph{subadditivity}; even for disjoint
$A,B$ it need not satisfy $\lambda^*(A\cup B)=\lambda^*(A)+\lambda^*(B)$.
If it were a countably additive measure on all $\mathcal P(\R)$,
its translation invariance and correct interval lengths would
contradict the Vitali-set argument from Module 3. The remedy is to
\emph{restrict the domain} to sets on which splitting is additive.

\key{The covering construction assigns an outer size to every set, but
the full power set is too large a domain for a translation-invariant,
countably additive extension of interval length.}

\section{Carath\'eodory's Criterion}
\subsection{Measurability by additive splitting}
Let $\mu^*$ be an outer measure on $\Omega$. A set $A\subseteq\Omega$
is \emph{Carath\'eodory measurable} if for \emph{every} $E\subseteq\Omega$,
\[
\boxed{\mu^*(E)=\mu^*(E\cap A)+\mu^*(E\setminus A).}
\]
The inequality ``$\le$'' follows automatically from subadditivity;
the criterion demands the reverse inequality as well. Thus a measurable $A$ splits any test set without increasing outer size;
testing only $E=A$ would be vacuous.

\textbf{Null sets are measurable.} If $\mu^*(A)=0$, then
$\mu^*(E\cap A)=0$ and $\mu^*(E\setminus A)\le\mu^*(E)$.
Together with subadditivity this proves the criterion. Every subset
of an outer-null set is likewise measurable.

\subsection{Carath\'eodory's theorem}
\textbf{Theorem.} Let $\mu^*$ be any outer measure. Then
\[
\boxed{\begin{aligned}
\mathcal M&=\{A\subseteq\Omega:A\text{ satisfies the splitting criterion}\}
\\[-2pt]
&\text{is a sigma-algebra, and}\qquad
\mu=\mu^*|_{\mathcal M}\text{ is a measure.}
\end{aligned}}
\]
\emph{Proof sketch.} The empty set is measurable, and complementation
just exchanges the two pieces of a split. Repeated splitting establishes
finite unions. For pairwise disjoint measurable $A_n$, repeated
application of the criterion, followed by countable subadditivity,
shows that $A=\bigcup_nA_n$ is measurable and
\[
\mu^*(A)=\sum_n\mu^*(A_n).
\]
Disjointification extends this to arbitrary countable unions.

\key{Carath\'eodory's test converts a universally defined, merely
subadditive outer measure into a genuinely countably additive measure
on a sigma-algebra.}

\section{Lebesgue Measure}
\subsection{Definition and main properties}
Apply Carath\'eodory's theorem to Lebesgue outer measure $\lambda^*$:
\[
\boxed{\begin{gathered}
\Lcal=\{A\subseteq\R:A\text{ is Carath\'eodory measurable for }\lambda^*\},\\
\lambda=\lambda^*|_{\Lcal},\qquad(\R,\Lcal,\lambda).
\end{gathered}}
\]
This is \emph{Lebesgue measure}, with the following properties:
\begin{itemize}
\item \textbf{Length:} $\lambda([a,b])=b-a$; hence $\lambda(\R)=\infty$.
\item \textbf{Translation invariance:} $A\in\Lcal\Rightarrow
A+t\in\Lcal$ and $\lambda(A+t)=\lambda(A)$.
\item \textbf{Countable additivity:} $\lambda(\bigsqcup_n A_n)=\sum_n\lambda(A_n)$.
\item \textbf{Completeness:} every subset of a Lebesgue null set is Lebesgue measurable and null.
\end{itemize}
Measure-theoretic completeness is \emph{not} Cauchy completeness from
metric-space theory.

\subsection{Borel sets and Lebesgue measure}
Every open subset of $\R$ is Lebesgue measurable (proved by splitting
interval covers, first for intervals and then for open sets). Since
$\Lcal$ is a sigma-algebra containing the usual open sets,
\[
\boxed{\B(\R)=\sigma(\mathcal T_{\mathrm{usual}})\subseteq\Lcal.}
\]
The Borel sigma-algebra is generated by topology; Lebesgue measurability
is determined by additive splitting of outer measure.

\subsection{Borel versus Lebesgue measurability}
Both inclusions are strict:
\[
\boxed{\B(\R)\subsetneq\Lcal\subsetneq\mathcal P(\R).}
\]
\textbf{Cantor set (middle-thirds construction).} Start with $C_0=[0,1]$;
at each step remove the open middle third of every remaining interval:
\[
C_1=[0,\tfrac13]\cup[\tfrac23,1],\qquad
C=\bigcap_{n=0}^{\infty}C_n.
\]
Each $C_n$ is a union of $2^n$ closed intervals of length $3^{-n}$.
Thus $C$ is closed (hence compact and Borel), and
\[
\boxed{0\le\lambda(C)\le\lambda(C_n)=(\tfrac23)^n\to0,
\qquad\lambda(C)=0.}
\]
Yet $C$ is \emph{uncountable}: its elements are exactly the numbers with
ternary expansions using only $0$ and $2$, so $|C|=2^{\aleph_0}$.
This is an uncountable compact set of measure zero.

\textbf{Why the inclusions are strict.} Every $S\subseteq C$ is Lebesgue
measurable by completeness. But
$|\mathcal P(C)|=2^{2^{\aleph_0}}>2^{\aleph_0}=|\B(\R)|$,
so some such $S$ cannot be Borel. A Vitali set is not Lebesgue
measurable, proving the second strict inclusion.

\key{The construction is \emph{interval length} $\to$ \emph{outer measure}
$\to$ \emph{Carath\'eodory measurability} $\to$ \emph{Lebesgue measure}.
The restriction of the domain restores countable additivity.}

\end{document}
